\section{Description du problème}
\label{sec:probleme}

Cette section présente le problème de planification de maintenance ferroviaire. Le problème consiste à planifier les opérations de maintenance sur un horizon de 24 heures en respectant les contraintes de ressources et les fenêtres temporelles des trains.

\subsection{Paramètres d'entrée}

Une instance du problème contient les éléments suivants :

\textbf{Trains :} Chaque train $t$ a un temps d'arrivée $a_t$ et un temps de départ $d_t$ (en minutes). Les maintenances doivent être effectuées pendant que le train est présent à la gare.

\textbf{Maintenances :} Chaque maintenance $m$ est associée à un train, a une priorité (1 = urgente, 2 = préventive) et une date butoir idéale $\delta_m$. L'objectif est de minimiser les retards par rapport à cette date butoir, en donnant plus d'importance aux maintenances urgentes.

\textbf{Opérations :} Chaque maintenance est composée de plusieurs opérations. Chaque opération $o$ a une durée $dur_o$, peut avoir un prédécesseur (contrainte de précédence), et nécessite une voie spécifique et des ressources (techniciens et équipements).

\textbf{Ressources :} 
\begin{itemize}
    \item \textbf{Techniciens :} Il y a plusieurs types de techniciens. Chaque type $k$ a un nombre limité $N_k$ de techniciens disponibles, et chaque technicien peut travailler au maximum $T_{max,k}$ minutes.
    \item \textbf{Équipements :} Il y a plusieurs types d'équipements. Chaque type $e$ a une capacité simultanée maximale $C_e$ (nombre d'opérations qui peuvent utiliser cet équipement en même temps).
    \item \textbf{Voies :} Les opérations doivent être assignées à des voies. Certains équipements ne sont disponibles que sur certaines voies.
\end{itemize}

\textbf{Setup :} Il y a des temps de setup entre opérations ($\sigma_{op}$ minutes) et entre trains différents ($\sigma_{train}$ minutes, plus long car il faut changer de train).

\subsection{Variables de décision}

Pour résoudre le problème, nous devons décider :
\begin{itemize}
    \item Le temps de début $s_o$ et de fin $f_o$ de chaque opération $o$
    \item Le temps de début $s_m$ et de fin $f_m$ de chaque maintenance $m$
    \item La voie $x_o$ assignée à chaque opération $o$
\end{itemize}

\subsection{Contraintes}

Une solution est valide si elle respecte toutes les contraintes suivantes :

\textbf{Cohérence temporelle :} La fin d'une opération est égale à son début plus sa durée : $f_o = s_o + dur_o$. Le début d'une maintenance est le minimum des débuts de ses opérations, et sa fin est le maximum des fins de ses opérations.

\textbf{Fenêtres temporelles :} Toutes les opérations d'une maintenance doivent être effectuées pendant que le train est présent : $a_t \leq s_m \leq f_m \leq d_t$ où $t$ est le train associé à la maintenance $m$.

\textbf{Précédence :} Si une opération a un prédécesseur, elle doit commencer après la fin du prédécesseur plus le temps de setup : $s_o \geq f_{pred(o)} + \sigma_{op}$.

\textbf{Non-chevauchement sur les voies :} Deux opérations ne peuvent pas utiliser la même voie en même temps. Si elles sont du même train, il faut un setup de $\sigma_{op}$ minutes entre elles. Si elles sont de trains différents, il faut un setup de $\sigma_{train}$ minutes.

\textbf{Capacité des ressources :} À tout moment, le nombre de techniciens de chaque type utilisés ne doit pas dépasser le nombre disponible. De même, le nombre d'équipements de chaque type utilisés simultanément ne doit pas dépasser la capacité maximale.

\textbf{Compatibilité voie-équipement :} Si une opération nécessite un équipement, elle doit être assignée à une voie où cet équipement est disponible.

\textbf{Limite de charge de travail :} La charge totale de travail (somme des durées × nombre de techniciens) pour chaque type de technicien ne doit pas dépasser la capacité totale disponible.

\subsection{Fonction objectif}

L'objectif est de minimiser le retard pondéré total : $\min \sum_{m} p_m \cdot \max(0, f_m - \delta_m)$ où $p_m$ est la priorité de la maintenance $m$ et $\delta_m$ est sa date butoir. Les maintenances urgentes (priorité 1) ont un poids plus élevé que les préventives (priorité 2).
